% !TeX root = ../偏微分方程笔记.tex
% 由主文档通过 \input 载入；请编译主文档。
\setcounter{chapter}{3} % 保留原书章号。
\chapter{代表习题的解答与提示}
\label{app:worked-exercises}

本附录选取各章一个代表问题，补充关键推理与核查步骤，不替代独立练习。计算题应先判断问题结构，再尝试推导；证明题应先写出要使用的假设，避免从一个正确公式跳到一个缺少条件的结论。

\section{基本概念：平移为何保持误差}
对应第\ref{chap:concepts}章。设两组有界初值为 $g_1,g_2$，输运解为 $u_j(x,t)=g_j(x-at)$。因为映射 $x\mapsto x-at$ 是全轴到自身的双射，
\[
 \sup_{x\in\mathbb R}|u_1(x,t)-u_2(x,t)|
 =\sup_{y\in\mathbb R}|g_1(y)-g_2(y)|.
\]
因此误差既不放大也不缩小。在有限空间区间上，平移可能把追踪点带出区间，此时还要计入流入边值，不能不加说明地使用上述全轴等式。若初值只作为 $L^\infty$ 等价类给定，则把上确界换成本性上确界即可。

\section{物理模型：有热源时是否存在稳态}
对应第\ref{chap:modeling}章。对 $u_t-\kappa u_{xx}=q$、$q>0$，若两端绝热，则积分给出 $d\bar u/dt=q$，其中 $\bar u=\ell^{-1}\int_0^\ell u\,dx$。于是 $\bar u(t)=\bar g+qt$，不存在温度有界的稳态。若两端保持零温，稳态是 $v=q x(\ell-x)/(2\kappa)$。此时左端流出量为 $\kappa v'(0)=q\ell/2$，右端流出量为 $-\kappa v'(\ell)=q\ell/2$，两者之和恰好平衡总源量 $q\ell$。这里采用已除以热容量的归一化通量；恢复物理单位时需乘回相应常数。

\section{特征线：从竖直初始线恢复参数}
对应第\ref{chap:characteristics}章。在 $x=1$ 上给 $u(1,y)=g(y)$，求 $xu_x+yu_y=u$。从 $(1,r)$ 出发的特征为 $x=e^s$、$y=re^s$，沿线有 $z'=z$、$z(0)=g(r)$，所以 $z=e^sg(r)$。在 $x>0$ 可反解 $s=\log x$、$r=y/x$，得到
\[
 u(x,y)=xg(y/x).
\]
直接求导得 $u_x=g(r)-rg'(r)$、$u_y=g'(r)$，所以 $xu_x+yu_y=xg(r)=u$。$x=1$ 的非特征行列式为 $-1$。若只在初始线的一段给数据，则公式只确定由该线段发出的特征区域；初始线不能自动决定 $x<0$ 中的解。

\section{方程分类：特里科米方程的特征线}
对应第\ref{chap:classification}章。在 $y<0$，方程 $yu_{xx}+u_{yy}=0$ 的特征关系为 $y(dy)^2+(dx)^2=0$。写成 $dx/dy=\pm\sqrt{-y}$ 并积分，得到两族
\[
 x+\frac23(-y)^{3/2}=C_1,\qquad
 x-\frac23(-y)^{3/2}=C_2.
\]
把左侧作为新坐标时，雅可比为 $2\sqrt{-y}\ne0$，但在 $y=0$ 趋于零。这解释了为什么双曲区内的坐标不能作为穿过退化线的同一组非退化坐标。处处椭圆却不一致椭圆的另一个例子是 $\mathbb R^2$ 上 $M(x)=\operatorname{diag}(1,(1+|x|^2)^{-1})$，其最小特征值为正，却没有统一正下界。

\section{傅里叶级数：由系数证明庞加莱不等式}
对应第\ref{chap:fourier-series}章。先取 $v\in C^1([0,\ell])$ 且 $v(0)=v(\ell)=0$，设其正弦系数为 $b_k$。分部积分说明 $v'$ 的第 $k$ 个余弦系数为 $(k\pi/\ell)b_k$，常数系数为零。两次使用帕塞瓦尔等式，得到
\[
 \int_0^\ell v^2dx=\frac\ell2\sum_{k\ge1}b_k^2,\qquad
 \int_0^\ell(v')^2dx=\frac\ell2\sum_{k\ge1}\left(\frac{k\pi}{\ell}\right)^2b_k^2
 \ge\left(\frac\pi\ell\right)^2\int_0^\ell v^2dx.
\]
等号当且仅当只有第一模态可能非零，即 $v$ 为 $\sin(\pi x/\ell)$ 的常数倍。由 $H_0^1$ 中的稠密性可推广到能量空间。这样既证明了不等式，也找到了最优常数及其达到者。

\section{分离变量：恒定模态源的响应}
对应第\ref{chap:separation}章。求零初值、零端点的 $u_t-\kappa u_{xx}=\sin(\pi x/\ell)$。令 $\lambda_1=(\pi/\ell)^2$，代入 $u=a(t)\sin(\pi x/\ell)$ 得 $a'+\kappa\lambda_1a=1$、$a(0)=0$。积分因子给出
\[
 u(x,t)=\frac{1-e^{-\kappa\lambda_1t}}{\kappa\lambda_1}\sin\frac{\pi x}{\ell}.
\]
直接求导时，时间项贡献 $e^{-\kappa\lambda_1t}$，负的扩散项贡献 $1-e^{-\kappa\lambda_1t}$，两者之和为单位振幅的源。稳态是 $(\kappa\lambda_1)^{-1}\sin(\pi x/\ell)$，说明长期响应由源与边界散热共同决定。

\section{一维波动：波经过后能否留下位移}
对应第\ref{chap:wave-one}章。令 $g=0$，并设 $h$ 光滑且支集包含在 $[a,b]$ 内。对固定 $x$，只要 $ct>\max\{|x-a|,|x-b|\}$，达朗贝尔积分已覆盖整个支集，因此
\[
 u(x,t)=\frac1{2c}\int_a^b h(s)ds.
\]
若积分不为零，固定位置会留下常数位移；若积分为零，信号经过后恢复零位移。这不违背有限传播，因为任意有限时刻解仍在有界传播区间之外为零；结论是先固定位置、再让时间足够大。对常值源 $f=1$、零初值的全轴问题，杜阿梅尔公式则给出 $u=t^2/2$，直接代入有 $u_{tt}=1$、$u_{xx}=0$。

\section{多维波动：阻尼和外力怎样改变能量}
对应第\ref{chap:wave-multi}章。对足够光滑的 $u_{tt}+\gamma u_t-c^2\Delta u=f$，在零边界能量通量条件下乘以 $u_t$ 并积分，得到
\[
 E'(t)+\gamma\|u_t\|_2^2=(f,u_t)_{L^2}.
\]
若 $\gamma>0$，杨氏不等式进一步给出 $E'+(\gamma/2)\|u_t\|_2^2\le\|f\|_2^2/(2\gamma)$。阻尼消耗能量，外力做功则可能为正也可能为负。没有源时只得到单调耗散；若要证明总能量指数衰减，还需额外论证速度耗散怎样间接控制位移能量。

\section{热方程：定量描述逆向误差放大}
对应第\ref{chap:heat}章。设 $\phi_m$ 为 $L^2$ 规范的正弦模态，$\lambda_m=(m\pi/\ell)^2$。在终止时刻 $T$ 加入误差 $\delta\phi_m$，其中 $\delta>0$，则与之对应的初值误差为 $\delta e^{\kappa\lambda_mT}\phi_m$。给定任意希望达到的误差放大量 $A>0$，只需取足够大的 $m$，使 $\delta e^{\kappa\lambda_mT}>A$。这个反例使用的终态都是确实可达的光滑函数，因此不适定性不能归咎于选择了不存在解的数据。

\section{调和函数：圆环中的径向解}
对应第\ref{chap:harmonic}章。对圆环 $a<r<b$、$0<a<b$，径向调和方程为 $(rU')'=0$，故 $U=A\log r+B$。给定 $U(a)=T_a$、$U(b)=T_b$，得到
\[
 U(r)=T_a+(T_b-T_a)\frac{\log(r/a)}{\log(b/a)}.
\]
它在两边界温度之间，符合最大值原理。如果区域包含原点而要求有界，对数项必须消失，因此整个圆盘中的有界径向调和函数只能为常数。原点条件改变允许的解族，不能把圆环公式直接用于实心圆盘。

\section{泊松方程：先检查总源与边界通量}
对应第\ref{chap:green}章。若 $-\Delta u=1$、$\partial_\nu u=0$，则相容条件要求 $|\Omega|=0$，与非空有界区域矛盾，故没有相应经典解或弱解。若在半径 $R$ 的球上改取常数法向导数 $g=-R/n$，由 $|B_R|=R|\partial B_R|/n$ 可知相容条件成立，且
\[
 u(x)=-\frac{|x|^2}{2n}+C
\]
直接满足方程与边界。若还要求零均值，利用球内 $|x|^2$ 的平均值为 $nR^2/(n+2)$，得到 $C=R^2/(2(n+2))$。边界积分的符号与常数歧义在这个例子中都可逐项核查。

\section{积分变换：平移与伸缩}
对应第\ref{chap:transforms}章。按本书一维变换约定，令 $y=x-a$，得到 $\widehat{f(\cdot-a)}(\xi)=e^{-ia\xi}\widehat f(\xi)$；令 $y=bx$、$b>0$，得到 $\widehat{f(b\cdot)}(\xi)=b^{-1}\widehat f(\xi/b)$。空间平移改变相位，空间压缩则拓宽频率尺度，同时产生雅可比因子。对 $f\in L^1$ 可直接换元证明；$L^2$ 情形通过稠密性延拓。高维同一伸缩的因子为 $b^{-n}$，不能沿用一维的 $b^{-1}$。

\section{索伯列夫空间：能量控制连续代表}
对应第\ref{chap:sobolev}章。对区间上的光滑函数和 $x>y$，微积分基本定理及柯西--施瓦茨给出
\[
 |v(x)-v(y)|=\left|\int_y^xv'(s)ds\right|
 \le\|v'\|_{L^2}|x-y|^{1/2}.
\]
对 $H^1$ 逼近序列应用同样估计，并结合端点或平均值控制，可得到一致收敛的连续代表，极限仍满足此不等式。这是一维能量函数具有连续代表的原因；高维不能把线段上的积分直接用全区域 $L^2$ 范数控制，所以不推出任意维数中的 $H^1$ 函数都连续。

\section{椭圆弱解：共振时的可解与不可解}
对应第\ref{chap:elliptic-weak}章。在 $(0,\ell)$ 上取零端点，考虑 $-u''-\lambda_1u=f$。若 $f=\sin(\pi x/\ell)$，将弱方程测试以同一正弦，左侧为零，右侧为 $\ell/2>0$，矛盾。若 $f=\sin(2\pi x/\ell)$，则
\[
 u(x)=\frac{\sin(2\pi x/\ell)}{\lambda_2-\lambda_1}
       +C\sin(\pi x/\ell),\qquad C\in\mathbb R,
\]
都是解。前一问题无解，后一问题有无穷多解，恰好对应源项是否与零空间正交。算子在共振参数处失去强制性，不能直接套用正定问题的唯一性结论。

\section{演化弱解：初值和源项同时扰动}
对应第\ref{chap:evolution-weak}章。记差解为 $w$、初值差为 $w_0$、源项差为 $r\in L^2(0,T;V^*)$。测试差方程以 $w$，并使用 $\|v\|_V=\|\nabla v\|_2$，得到
\[
 \frac12\frac d{dt}\|w\|_H^2+\kappa\|w\|_V^2
 \le\frac\kappa2\|w\|_V^2+\frac1{2\kappa}\|r\|_{V^*}^2.
\]
积分后即有 $\|w(t)\|_H^2+\kappa\int_0^t\|w\|_V^2ds\le\|w_0\|_H^2+\kappa^{-1}\int_0^t\|r\|_{V^*}^2ds$。这是一条含全部数据误差的稳定估计，也说明选择对偶范数怎样配合能量空间。

\section{守恒律：交通流中的逆向拥堵波}
对应第\ref{chap:conservation}章。通量 $F(\rho)=v_0\rho(1-\rho/\rho_{\max})$ 满足 $F''=-2v_0/\rho_{\max}<0$，故为凹通量。左右密度为 $\rho_L,\rho_R$ 的间断速度为
\[
 s=v_0\left(1-\frac{\rho_L+\rho_R}{\rho_{\max}}\right).
\]
当 $\rho_L<\rho_R$ 时，左侧特征比右侧快，形成进入型激波；这与凸伯格斯通量的状态次序相反。若 $\rho_L=\rho_{\max}/2$、$\rho_R=\rho_{\max}$，则 $s=-v_0/2$，拥堵边界向上游移动，而车辆速度仍非负。车辆的移动速度与密度扰动的传播速度是不同对象。

\section{非线性方法：全局利普希茨反应项的上界}
对应第\ref{chap:nonlinear}章。设 $F$ 的全局利普希茨常数为 $L\ge0$，则 $|F(r)|\le|F(0)|+L|r|$。温和解及热半群收缩性给出
\[
 \|u(t)\|_\infty\le\|g\|_\infty+t|F(0)|
       +L\int_0^t\|u(s)\|_\infty ds.
\]
若 $L>0$，积分形式的格朗沃尔估计给出 $\|u(t)\|_\infty\le e^{Lt}\|g\|_\infty+|F(0)|(e^{Lt}-1)/L$；若 $L=0$，则为 $\|g\|_\infty+t|F(0)|$。任意有限时间内该界有限，结合局部延拓准则得到整体温和解。整体存在不意味着时间趋于无穷时有统一界。

\section{数值方法：从放大因子识别不稳定}
对应第\ref{chap:numerics}章。在周期偶数节点网格上取交替符号模态 $U_j^m=A_m(-1)^j$，则显式热格式给出 $A_{m+1}=(1-4\mu)A_m$。当 $\mu=1/4$，该最高频率一步消失；当 $\mu=1/2$，它每步变号但振幅不变；当 $\mu>1/2$，其绝对振幅按几何级数增长。稳定边界上的不衰减仍可能影响实际精度，所以稳定不等于准确。在固定零边界网格上应使用离散正弦模态，放大因子为 $1-4\mu\sin^2(k\pi/(2J))$；随网格加密，最高频率给出同样的统一步长限制。

